%=============================================================================
% WAVE 4, ITERATION 10: LENSING PAPER CROSS-CHECK
% Agent: Opus 4.6 | Date: 20 March 2026
% Assignment: Lensing cross-check against MASTER\_FINDINGS\_CORPUS
%=============================================================================

\section*{W4i10: Lensing Cross-Check --- Top 5 Findings}

\subsection*{Executive Summary}

Cross-referenced all lensing-related equations and claims in
\texttt{section02\_formalism.tex}, \texttt{section04\_ppn.tex},
\texttt{section07\_galactic.tex}, \texttt{section12\_cosmology.tex},
\texttt{appendix\_B.tex}, \texttt{appendix\_I.tex}, and
\texttt{appendix\_J.tex} against the definitive
\texttt{MASTER\_FINDINGS\_CORPUS.md}.  No lensing-specific standalone paper
(``DFD\_Lensing\_Paper\_Overleaf.zip'' or equivalent) was found in the
repository; the lensing content is distributed across the unified review
sections listed above.

Five findings are ranked by severity: two weaknesses that referees
\emph{will} attack, two internal inconsistencies/gaps, and one opportunity
to strengthen the mu-shape test.

%=============================================================================
\subsection*{Finding 1 [WEAKNESS --- HIGHEST PRIORITY]:
Bullet Cluster Lensing Offset Is Under-Derived}

\paragraph{Problem.}
Section~7 (Table~\ref{tab:bullet} equivalent, lines~700--712 of
\texttt{section07\_galactic.tex}) presents the Bullet Cluster lensing
offset as ``DFD offset = 129~kpc (main) and 163~kpc (bullet
subcluster)'' matching observations at 83\% and 72\%.  However:

\begin{enumerate}
\item \textbf{No derivation is given.}  The claimed offsets appear as
  table entries without any supporting calculation.  No projected
  surface-density integral, no lens-model geometry, no specification of
  which baryonic mass profile was assumed.
\item \textbf{The mechanism sketch is qualitative only.}  The text states
  ``forces cancel at gas center $\to$ large $\Psi$'' and ``net effect
  shifts lensing peak toward galaxies'' but gives no equation for the
  effective convergence $\kappa_{\rm DFD}(R) = \Sigma_{\rm bar}(R) \cdot
  \Psi(R) / \Sigma_{\rm crit}$.
\item \textbf{72\% match for the subcluster is poor.}  A 28\% miss on
  a flagship system will be immediately flagged by referees familiar
  with Clowe et al.\ (2006), Bradac et al.\ (2006), and Markevitch
  (2006).
\end{enumerate}

\paragraph{Referee objection.}
``The Bullet Cluster analysis is hand-waving.  Where is the projected
surface-density map?  Where is the convergence profile?  The 72\% match
for the subcluster is a 28\% discrepancy on the single most important
cluster for dark matter claims.  Without a rigorous N-body or
hydrodynamic simulation, this is not a resolution.''

\paragraph{Recommended action.}
\begin{itemize}
\item Provide an explicit equation for $\kappa_{\rm DFD}(\mathbf{R})$
  based on $\Sigma_{\rm eff} = \Sigma_{\rm bar} \times \Psi(a/a_0)$
  integrated along the line of sight.
\item Specify the baryonic mass model (X-ray gas + BCG + ICL).
\item Run (or cite) a quantitative projected-mass reconstruction.
\item If 72\% cannot be improved, explicitly state the residual and
  discuss whether ram-pressure geometry or projection effects could
  close the gap.
\end{itemize}

%=============================================================================
\subsection*{Finding 2 [WEAKNESS --- HIGH PRIORITY]:
Cluster ``Resolution'' Depends Heavily on Two Ad Hoc Corrections}

\paragraph{Problem.}
Appendix~I (Table~\ref{tab:cluster-final} equivalent) claims
``Obs/DFD = $0.98 \pm 0.05$'' for all 16 clusters, but this result
requires \emph{two substantial corrections} applied simultaneously:

\begin{enumerate}
\item \textbf{Baryonic mass correction:} 25--45\% upward revision
  (WHIM + ICL + clumping + gas beyond $r_{500}$).
\item \textbf{Jensen averaging correction:} 25--46\% from subhalo
  convexity of $\Psi = 1/\mu$.
\end{enumerate}

The \emph{raw} Obs/DFD ratios are 1.47--2.14 for clusters ---
systematically off by factors of 1.5--2$\times$.  The corrections bring
these into agreement, but:

\begin{itemize}
\item Both corrections push in the \emph{same} direction (lowering
  Obs/DFD).  This is suspicious: two independent corrections that
  both happen to improve the fit.
\item The subhalo mass fractions ($f_{\rm sub} = 0.14$--0.27) are
  \emph{assumed}, not measured per cluster.
\item The baryonic correction factors (1.25--1.45) span a wide range
  and are assigned without per-cluster justification.
\end{itemize}

\paragraph{Referee objection.}
``Both your baryonic and Jensen corrections are tuned to bring the
clusters into agreement.  Without independent per-cluster measurements
of $f_{\rm sub}$ and missing baryonic mass, this is effectively two
free parameters disguised as `physically motivated corrections.'
The raw ratios of 1.5--2 are the honest numbers.''

\paragraph{Recommended action.}
\begin{itemize}
\item Cite specific per-cluster measurements of subhalo mass fractions
  from lensing surveys (e.g., Merten et al.\ 2011, Okabe et al.\ 2014).
\item State the baryonic correction as a range with per-cluster error
  bars, not point values.
\item Explicitly acknowledge that the raw ratios are 1.5--2$\times$
  and that the corrections are necessary conditions, not ancillary
  refinements.
\item Consider presenting both raw and corrected results as separate
  columns with equal prominence.
\end{itemize}

%=============================================================================
\subsection*{Finding 3 [INTERNAL GAP]:
$\psi$-Lensing $\Delta\psi = 0.30$ vs.\ Screen $\Delta\psi = 0.27$ at
$z = 1$}

\paragraph{Problem.}
Two different values of $\Delta\psi$ at $z \sim 1$ appear in the paper:

\begin{itemize}
\item \textbf{Appendix~J} (CMB peak location): $\Delta\psi = 0.30$
  required to shift $\ell_{\rm true} = 297 \to \ell_{\rm obs} = 220$.
  (This is $\psi_{\rm CMB} - \psi_{\rm here}$.)
\item \textbf{Section~12} (cosmological screen): $\Delta\psi_{\rm
  screen} \approx 0.27$ at $z \sim 1$ (stated explicitly at line~651
  and in the tomography figures).
\end{itemize}

These are \emph{different physical quantities}: the CMB value integrates
from $z = 0$ to $z = 1100$, while the screen value is at $z = 1$.
However, the paper does not explicitly state how $\Delta\psi(z)$ evolves
from 0.27 at $z = 1$ to 0.30 at $z = 1100$.  The near-equality is
suggestive that most of the screening occurs at $z < 1$, but this is
never stated.

\paragraph{Referee objection.}
``What is $\Delta\psi(z)$?  You quote 0.27 at $z = 1$ and 0.30 at the
CMB.  Does $\Delta\psi$ saturate?  What is the functional form?  Without
an explicit $\Delta\psi(z)$, the CMB peak-location analysis and the
cosmological screen are disconnected claims, not a unified prediction.''

\paragraph{Recommended action.}
\begin{itemize}
\item Provide an explicit $\Delta\psi(z)$ curve or functional form,
  even if approximate.
\item Show that $\Delta\psi(z=1100) \approx 0.30$ is consistent with
  $\Delta\psi(z=1) \approx 0.27$ via the integral of the growth
  equation or the dust-branch evolution.
\item If the two values are essentially independent claims, state
  this explicitly.
\end{itemize}

%=============================================================================
\subsection*{Finding 4 [OPPORTUNITY]:
$\mu(x)$ Shape Test --- Galaxy-Galaxy Lensing Is the Tightest Constraint}

\paragraph{Current status.}
The corpus states that $\mu(x) = x/(1+x)$ is ``uniquely forced by
$S^3$ composition'' (Theorem in Appendix~N).  The paper calibrates
once on the RAR and freezes.  The question is: what observations pin
the \emph{shape} of $\mu(x)$ most tightly in the transition region
$x \sim 0.1$--$10$?

\paragraph{Analysis of constraining power.}
\begin{enumerate}
\item \textbf{SPARC RAR:} Constrains $\mu(x)$ in the range $x \sim
  0.01$--$10$.  This is the current calibration anchor.  Residual
  scatter $\sim 0.13$ dex leaves the transition shape mildly
  degenerate between simple and standard $\mu$.
\item \textbf{Galaxy-galaxy lensing (stacked):}  Probes $x \sim
  0.01$--$1$ via the effective surface density $\Sigma_{\rm eff}
  \propto \Sigma_{\rm bar}/\mu$.  Brouwer et al.\ (2021, A\&A) measured
  the weak-lensing RAR extending 2 decades below galaxy rotation curve
  data.  This is the \textbf{tightest available shape test} because:
  (a) it probes the deep-field regime $x \ll 1$ where simple and
  standard $\mu$ diverge most; (b) stacking over $\sim 10^5$ galaxies
  reduces statistical noise below the RAR scatter.  A 2025 ApJ paper
  (Mistele et al.) proposes disk-galaxy lensing number counts as a
  novel MOND test; Euclid data (first cosmology release October 2026)
  will be decisive.
\item \textbf{Cluster lensing:} Probes $x \sim 0.05$--$0.1$ but is
  confounded by the Jensen averaging correction and baryonic
  uncertainties (Finding~2).  Currently not a clean shape test.
\item \textbf{CMB lensing convergence:} Probes the integrated $\psi$
  along the line of sight.  Sensitive to $G_{\rm eff}(k)$ from
  Eq.~(section12, eq.\ G-eff) but degenerate with the background
  $\Delta\psi(z)$ evolution.
\end{enumerate}

\paragraph{Recommended protocol for strengthening.}
\begin{itemize}
\item \textbf{Tier 1 (immediate):} Confront DFD $\mu(x) = x/(1+x)$
  with the Brouwer et al.\ (2021) weak-lensing RAR data.  Compute
  predicted $\Delta\Sigma(R)$ profiles and compare with KiDS/GAMA
  measurements.  This test can discriminate simple vs.\ standard $\mu$
  at $> 3\sigma$ using existing data.
\item \textbf{Tier 2 (2026--2027):} Use Euclid galaxy-galaxy lensing
  catalogs to measure the RAR at $x < 0.01$.  DFD predicts
  $g_{\rm obs}/g_{\rm bar} = 1/\mu(x) \approx 1/x$ in the deep-field
  limit; any deviation from $1/x$ scaling at very low $x$ would
  constrain the $\mu$ shape.
\item \textbf{Tier 3 (2027+):} Disk galaxy lens number counts
  (Mistele et al.\ 2025, ApJ).  MOND-like theories predict
  $\sim 2\times$ more disk lenses than $\Lambda$CDM; this is a
  clean count test independent of profile modeling.
\end{itemize}

%=============================================================================
\subsection*{Finding 5 [WEAKNESS]:
$E_G$ Gravity Test and KiDS-Legacy $S_8$ Noted as ``Mild Tension''
but Not Quantified}

\paragraph{Problem.}
Section~12 (lines 819--832) acknowledges two counter-evidence items:

\begin{enumerate}
\item \textbf{$E_G$ from ACT DR6 + BOSS:} ``consistent with
  $\Lambda$CDM/GR and largely scale-independent.''  The paper says DFD
  ``would expect small deviations at low $z$'' but gives no number.
\item \textbf{KiDS-Legacy $S_8$:} ``consistent with Planck $\Lambda$CDM.''
  The paper notes ``earlier KiDS analyses showed larger discrepancy.''
\end{enumerate}

The corpus (Section~3, Prediction~1) states that DFD predicts
$S_8^{\rm DFD} \sim 0.77$--$0.79$, and that the S$_8$ tension between
WL and Planck \emph{is} a DFD prediction.  If KiDS-Legacy now agrees
with Planck, that \emph{removes} DFD's claimed success.

\paragraph{Referee objection.}
``You claim the S$_8$ tension is a DFD prediction.  But KiDS-Legacy
(Li et al.\ 2023) now finds $S_8 = 0.815 \pm 0.016$, consistent with
Planck.  Does this falsify your prediction?  And your $E_G$ discussion
is purely qualitative --- what value does DFD predict for $E_G$?''

\paragraph{Quantitative assessment.}
\begin{itemize}
\item KiDS-Legacy $S_8 = 0.815 \pm 0.016$:  DFD predicts 0.77--0.79.
  This is a $1.5$--$2.8\sigma$ discrepancy depending on the DFD value.
  Not falsifying, but the earlier ``DFD predicts the tension'' claim
  is weakened.
\item The DFD prediction is scale-dependent: $S_8 \sim 0.73$--0.76
  at large scales ($> 100'$) and $\sim 0.81$--0.83 at small scales
  ($< 10'$).  KiDS-Legacy reports a single effective $S_8$.  If the
  scale-dependent prediction is correct, the KiDS number should
  depend on angular-scale cuts --- this is testable.
\item $E_G$: DFD's $G_{\rm eff}$ (Eq.~\eqref{eq:G-eff} in Section~12)
  predicts a direction-dependent correction.  The isotropically
  averaged $E_G$ at low $z$ should be $\sim 5$--10\% below GR.
  This should be stated explicitly.
\end{itemize}

\paragraph{Recommended action.}
\begin{itemize}
\item Compute DFD's predicted $E_G(z, \ell)$ using the $G_{\rm eff}$
  formula and compare with ACT DR6 values.
\item Update the $S_8$ discussion to acknowledge KiDS-Legacy and
  reframe: ``DFD predicts scale-dependent $S_8$; the single-number
  $S_8$ from KiDS-Legacy is intermediate, as expected if both
  high-$\ell$ and low-$\ell$ bins are averaged.''
\item Identify what angular-scale cuts in KiDS-Legacy or Euclid
  would reveal the predicted scale dependence.
\end{itemize}

%=============================================================================
\subsection*{Summary: Anticipated Referee Objections}

\begin{enumerate}
\item ``The Bullet Cluster analysis is qualitative and the 72\% match is
  poor.  Where is the convergence map?'' (Finding~1)
\item ``Your cluster resolution requires two ad hoc corrections totaling
  50--90\%.  The raw ratios are 1.5--2$\times$.'' (Finding~2)
\item ``What is $\Delta\psi(z)$?  The CMB and cosmological values appear
  inconsistent.'' (Finding~3)
\item ``KiDS-Legacy $S_8$ no longer shows tension with Planck.  Does this
  falsify your prediction?'' (Finding~5)
\item ``How does DFD explain the Skordis-Zlosnik (2021) result that
  relativistic MOND requires a vector field for lensing?  DFD has only
  a scalar --- is the lensing derivation complete?'' (This is the
  deepest theoretical question: DFD's $\gamma = 1$ from the
  exponential metric $n = e^\psi$ does give correct 1PN deflection,
  but the equivalence breaks down at cluster scales where $\psi$ is
  not small.  The $n = e^\psi$ form may need explicit verification
  beyond weak-field for cluster lensing convergence.)
\end{enumerate}

\subsection*{Recent Literature Notes (2025--2026)}

\begin{itemize}
\item Mistele et al.\ (2025, ApJ):  Disk galaxy gravitational lensing
  as a novel MOND test.  MOND predicts $\sim 2\times$ more disk lenses
  than CDM.  Testable with Euclid/Rubin.  Directly relevant to
  DFD $\mu$-shape.
\item Brouwer et al.\ (2021, A\&A):  Weak lensing RAR extending 2
  decades below rotation curve data.  The tightest current $\mu$-shape
  constraint.  DFD should confront this data directly.
\item Skordis \& Zlosnik (2021, PRL):  First relativistic MOND passing
  CMB + lensing.  Uses scalar + vector fields.  DFD's scalar-only
  approach should explicitly address why the vector field is not needed.
\item Euclid first cosmology data release: October 2026.  Galaxy-galaxy
  lensing, cosmic shear, and cluster lensing will be decisive for
  $\mu$-shape, $S_8$ scale dependence, and $E_G$.
\item ACT DR6 + BOSS $E_G$: Currently consistent with GR.  DFD
  should compute its predicted $E_G$ value.
\item IJFMR 2026 paper on Bullet Cluster without dark matter:
  Independent analysis worth cross-checking against DFD's treatment.
\end{itemize}

%=============================================================================
\subsection*{Cross-Check: Equations Verified Against Corpus}

\begin{tabular}{lcc}
\toprule
Equation / Claim & Status & Notes \\
\midrule
$\delta\theta = 4GM/(c^2 b)$ (1PN deflection) & CONSISTENT & Matches $\gamma = 1$ \\
$\alpha^{(2)}$: $15\pi/16$ coefficient (2PN) & CONSISTENT & Matches GR exactly \\
$\mu(x) = x/(1+x)$ (all lensing sections) & CONSISTENT & Uniquely derived (S$^3$) \\
$a_0 = 2\sqrt{\alpha}\,cH_0 = 1.12 \times 10^{-10}$ & CONSISTENT & Corrected at W3i7 \\
Shadow ratio $2e/(3\sqrt{3}) = 1.046$ & CONSISTENT & Confirmed to $10^{-17}$ \\
$\Delta\psi_{\rm screen} \approx 0.27$ at $z=1$ & INTERNALLY CONSISTENT & But see Finding~3 \\
$\Delta\psi_{\rm CMB} = 0.30$ & SEE FINDING 3 & Gap vs.\ 0.27 unexplained \\
Bullet offset 129/163~kpc & NO DERIVATION & See Finding~1 \\
Cluster Obs/DFD = $0.98 \pm 0.05$ (corrected) & METHODOLOGY CONCERN & See Finding~2 \\
$S_8 \sim 0.77$--$0.79$ & UNDER PRESSURE & KiDS-Legacy at 0.815 \\
$G_{\rm eff}$ (Eq.~G-eff, Section~12) & CONSISTENT & Properly derived \\
DDR $\eta = 1$ (Etherington) & CONSISTENT & Correctly maintained \\
\bottomrule
\end{tabular}

\vspace{1em}
\noindent\textit{Filed 20 March 2026.  Agent: Opus 4.6, Wave 4 Iteration 10.}

\end{document}
